% p1-01-introduction

\section{P1 §1. Introduction}

1. Introduction

    Figure (fig-p1-open). Opening triptych --- Seed Identity / Vesica Axis / Sprout. The trilogy opens
                                 from one identity into three papers.

  1.1 The Open Problem of Dimensionless Constants
  The numerical values of the fundamental dimensionless constants of nature --- the
  inverse fine-structure constant alpha-1 = 137.035 999 177(21) (NIST CODATA 2022),
  the strong coupling alphas(mZ) = 0.1180(9) (PDG 2024), the electroweak mixing angle
  sin2theta(MZ) = 0.23129(4) in the MS scheme (PDG 2024), and the full array of mixing
  angles, mass ratios, and cosmological parameters --- remain empirical inputs to the
  Standard Model. No known symmetry principle or dynamical mechanism within
  quantum field theory determines their specific values; each must be measured and
  inserted by hand. The question of whether these values carry any deeper algebraic or
  structural character is thus not closed by the Standard Model itself, and constitutes a
  legitimate open question at the intersection of mathematical physics and information
  theory (PDG 2024).

  1.2 A Precisely Bounded Reformulation
  A sharply bounded, operationally well-posed version of this question avoids dynamical
  claims entirely: Do the observed numerical values of dimensionless constants exhibit
  statistically non-random compressibility under a restricted symbolic language? This is a

  question about data compression, not about physics. It belongs to the domain of
  symbolic regression, algorithmic information theory, and statistical model selection
  under structured hypothesis classes --- the same theoretical territory inhabited by the
  Minimum Description Length principle (Barron, Rissanen \& Yu, IEEE Trans. Inf.
  Theory, 1998; Grünwald, MIT Press, 2007) and Chaitin's Algorithmic Information
  Theory (Chaitin, Cambridge University Press, 1987).

  This paper formalizes the question within a precisely defined statistical framework. A
  symbolic grammar G$\varphi$ is generated by the basis {$\varphi$, pi, e, 3} with integer exponents
  bounded by an ell1 complexity constraint C. The grammar induces a finite hypothesis
  class HC of candidate expressions. Physical constants form a frozen target dataset T,
  and the analysis asks whether the projection of T onto HC yields lower residuals than
  those expected from statistically appropriate null models.

  1.3 What This Paper Does Not Claim
  The following statements are explicitly disavowed and must be kept in view throughout:

  - No physical derivation. The grammar G$\varphi$ does not derive physical constants from
  symmetry principles, dynamics, or first principles. The encoding map T $\rightarrow$ S(C) is a
  constrained quantization onto a discrete algebraic manifold, not a generative physical
  law.
  - No ontological identification. The $\varphi$-basis is not claimed to carry physical primacy or to
  reflect any symmetry of nature. It is one instance of an algebraically complete basis
  with tractable combinatorial geometry.
  - No unique basis. The analysis would be qualitatively analogous for any other
  algebraically complete basis. The specific choice of {$\varphi$, pi, e, 3} is motivated by
  algebraic tractability and the Lucas ring Z[$\varphi$], as discussed in Sections 2.4 and 3.3.
  - No claim from stale proof counts. The PhD monograph program has ongoing formal
  verification work; no Coq-verified claims are made in this paper. Proof obligations are
  cited with their current honest-admission status.

  1.4 Primary Objectives
  The primary objectives of this paper are:

  1. To characterize the density of low-error approximations within HC as a function of
  complexity level C and tolerance $\varepsilon$.
  2. To assess whether this density exceeds the null expectation under three controlled
  baseline models, with full multiple-testing correction.
  3. To compute a well-defined MDL compression score that accounts for the full
  description cost of the symbolic model class.
  4. To introduce the Pellis hierarchical expansion and evaluate its contribution to
  compressibility under the MDL framework.
  5. To specify falsifiable predictions and out-of-sample tests that can in principle refute
  the compressibility hypothesis.
  6. To provide explicit negative controls, a falsification ledger, and a reviewer-risk
  assessment.

  1.5 Paper Organization
  Section 2 reviews related literature and positions the work. Section 3 defines the
  symbolic grammar and its combinatorial and algebraic structure in detail. Section 4
  specifies the target dataset and the pre-registration freeze protocol. Section 5 defines
  the null models. Section 6 addresses multiple testing corrections. Section 7 develops
  the MDL and Bayesian framework. Section 8 reports representative empirical results.
  Section 9 introduces the Pellis hierarchical expansion. Section 10 defines out-of-sample
  tests, a falsification ledger, and negative controls. Section 11 provides a reviewer-risk
  assessment. Section 12 discusses the relationship to the Trinity $S^{3}$AI / Flos Aureus PhD
  monograph, including the Stergios collaboration note. Section 13 discusses limitations.
  Section 14 concludes. Appendices provide the reproducibility protocol checklist,
  algebraic properties of the basis, phase-transition structure, and a mapping from the
  PhD monograph to Paper 1.

\subsection{References}

- Vasilev, D. \textit{Flos Aureus Monograph: Trinity Constants}. DOI \href{https://zenodo.org/doi/10.5281/zenodo.19227877}{10.5281/zenodo.19227877}.
- Pellis, S. "The Fine-Structure Constant and the Golden Ratio." \href{https://vixra.org/pdf/2110.0117v4.pdf}{viXra 2110.0117 v4}.
- Sherbon, M. A. "Fundamental Nature of the Fine-Structure Constant." \href{https://hal.science/hal-02341850/document}{HAL hal-02341850}.
